\documentclass[../../main.tex]{subfiles}
\begin{document}

\begin{center}
  {\fontsize{48pt}{58pt}\selectfont \textbf{1982}\normalsize\textbf{年度}}
\end{center}

\begin{table}[htb]
  \centering
  \caption{出題テーマと難易度}
  \begin{tabular}{|c|p{9cm}|c|}
    \hline
    \textbf{問題番号} & \textbf{テーマ} & \textbf{難易度} \\
    \hline
    第1問 & 円の配置と極限 & \\
    \hline
    第2問 & ベクトルと格子点 & \\
    \hline
    第3問 & 4次関数と接線の本数 & \\
    \hline
    第4問 & 楕円の回転と面積 & \\
    \hline
    第5問 & 確率 & \\
    \hline
  \end{tabular}
\end{table}

\begin{table}[htb]
  \centering
  \caption{解答}
  \renewcommand{\arraystretch}{1.4}
  \begin{tabular}{|c|c|p{9cm}|}
    \hline
    \textbf{問題番号} & \textbf{小問} & \textbf{答え} \\
    \hline
    第1問 & - & $\pi$ \\
    \hline
    第2問 & (1) & 証明問題 \\
    \cline{2-3}
    & (2) & $\dfrac{\sqrt{3}}{4}$ \\
    \hline
    第3問 & - & 図示（$D=\{(\alpha,\beta)\mid -1<\alpha<1,\ \alpha^4-6\alpha^2<\beta<3-8|\alpha|\}$，境界含まず） \\
    \hline
    第4問 & - & $\left(\dfrac{3}{2}+\dfrac{2\sqrt{3}}{3}\right)\pi$ \\
    \hline
    第5問 & - & $\dfrac{4}{27}$ \\
    \hline
  \end{tabular}
\end{table}

\paragraph{試験形態}
5問．試験時間は要確認．

\end{document}
